% verification.tex — 八、模型的检验与灵敏度分析（2026-08-15, 数字均来自测试记录）

\subsection{求解与实现的一致性检验}

为保证论文结果与程序实现完全一致，对全部模型代码执行了 26 项自动化一致性测试，
覆盖五个方面，全部通过：

\begin{enumerate}[label=（\arabic*）]
\item 维度检查（9 项）：各中间矩阵与结果表的规模均与设计一致，如品类日销量矩阵
      $6\times1085$、单品日销量矩阵 $251\times1085$、补货量与定价矩阵 $6\times7$、
      上架单品数落在 $[27,33]$ 区间内；
\item 特殊值与手算用例（5 项）：斯皮尔曼秩相关的手写实现与内置函数在含并列秩的
      样例上交叉验证一致（误差小于 $10^{-12}$）；三次多项式拟合在无噪声样本上
      精确恢复系数（误差小于 $10^{-6}$）；线性规划与 0-1 规划的小规模手算用例
      均得到手推最优解（补货 $(5,5)$、利润 25；选高价档、利润 15）；
\item 约束满足性（7 项）：在两个问题的全数据结果上逐条核验需求上限、价格区间、
      空间容量、最小陈列量与二元性约束，均在容差 $10^{-6}$ 内满足；
\item 稳定性（2 项）：对确定性求解流程重复运行 10 次，利润的标准差小于 $10^{-3}$，
      相对标准差小于 1\%；
\item 交叉验证（3 项）：利润按目标函数公式独立重算与求解器输出一致（误差小于
      $10^{-9}$）；空间容量放宽时利润单调不降；方案 B 在全数据上可行。
\end{enumerate}

\subsection{与历史数据的对比检验}

以方案 A 的口径回算 2020 至 2022 年 7 月 1 日至 7 日的利润（图12），2023 年
预测方案逐日利润均高于 2021、2022 年同期。2020 年同期利润偏高与该年疫情解封后
消费回补一致，可解释。该对比说明方案相对近两年经营水平确有提升，而非仅在模型
内部自洽。

\subsection{参数灵敏度与稳健性}

对关键假设参数做扰动检验：定价整体缩放 90\% 至 110\% 时（表6），七日利润在
4.58 千元至 8.89 千元之间变化，定价每变动 5\% 利润约同向变动 16\%，方向与幅度
均符合"价格是利润第一杠杆"的直觉，模型对定价参数的响应合理；空间容量在 0.8 至
1.2 倍区间扫描时，利润最大降幅 3.9\%，容量不低于基准值时利润不变，说明容量假设
不紧，模型对容量估计稳健；问题三的品类需求满足率下限从 0 提高到 50\% 时，最优解
不变，说明利润最优解天然满足品类均衡要求。综合三类扰动，模型的结论对参数假设
具有可接受的稳健性。

\subsection{检验小结}

上述检验的检验对象、指标与结论均可追溯（支撑材料含测试脚本与日志）。模型存在的
局限已在相应章节注明：茄类与辣椒类三次拟合的 $R^2$ 低于 0.1，其需求估计的置信
度有限；单品需求采用线性近似，样本不足单品依赖份额缩放。这些局限不影响问题二、
问题三主结论的方向，但提醒在应用方案时应结合后验销量数据滚动修正参数。
